% =========================================================
% CHAPTER 7: CONCLUSION AND FUTURE WORK
% =========================================================
% IMPORTANT: No sections or subsections allowed in Chapter 7.
% Paragraphs only, per Anna University guidelines.
% =========================================================

\chapter{CONCLUSION AND FUTURE WORK}
\label{chap:conclusion}

This project presented QueryPilot, an uncertainty-aware self-adaptive query
optimization framework for MySQL 8.0. The system was implemented as a Python
middleware that intercepts SQL queries, extracts structural features from
MySQL's EXPLAIN output, generates optimizer hint candidates, predicts their
performance using a three-model ensemble of Linear Regression, Random Forest,
and Gradient Boosting regressors, estimates prediction uncertainty through
ensemble disagreement, detects workload drift using the Page-Hinkley
sequential change-point test, and makes bounded-risk execution decisions
through the SAGE (Stability-Aware Guarded Execution) controller. The framework
was evaluated on the TPC-H benchmark at scale factor 1 under four controlled
workload scenarios using a methodologically rigorous interleaved execution
design that eliminates buffer pool warm-up bias. A real-time Streamlit
dashboard was built to visualize query performance, SAGE decisions, drift
detection events, and cross-scenario experiment comparisons.

The primary finding of this project is that conservative, uncertainty-aware
adaptation prevents the performance regressions that unconstrained ML-based
optimization introduces under dynamic workload conditions. QueryPilot + SAGE
achieved a regression rate of 0.0 percent in all four evaluated scenarios,
compared to regression rates of 10.0 to 14.0 percent for the ML-Only baseline
under the stable, sudden shift, and gradual shift scenarios. The per-query
overhead of QueryPilot's pipeline ranged from 3.21 to 5.86 milliseconds on
queries averaging 18 to 22 seconds in duration, representing less than 0.03
percent of total query time. These results confirm that a system that knows
when not to override the native optimizer is more reliable under workload
dynamics than one that always applies its best machine learning recommendation.

Several limitations of the current implementation must be acknowledged. First,
execution times are measured on the client side using Python's
time.perf\_counter, which includes result fetching and Python overhead in
addition to server-side query processing time. Server-side execution times
from MySQL's Performance Schema would provide a cleaner signal for the training
labels and are not used in the current version. Second, the candidate set is
small and restricted to three hint types supported by MySQL's optimizer hint
system. The available hint candidates (NO\_INDEX, HASH\_JOIN, NO\_BKA) do not
reliably outperform native execution on TPC-H with properly configured indexes,
which is why SAGE made zero adapt decisions across all scenarios. The
framework architecture is correct and the safety mechanisms work as intended,
but expanding the candidate pool to include more powerful hints would be
necessary for the system to demonstrate positive adaptations. Third, the model
retrains in batch after every ten observations, which may cause the model to
lag the current workload distribution under rapid shifts.

Future work on QueryPilot will focus on four directions. First, the candidate
generation strategy will be extended to include FORCE INDEX hints targeting
specific named indexes, JOIN ORDER hints for controlling multi-table join
order, and query rewriting strategies such as subquery-to-join transformations.
These hint types can provide genuine performance improvements on MySQL and
would give SAGE meaningful candidates to evaluate. Second, server-side
execution times from MySQL's Performance Schema will be used as training labels
to produce a cleaner signal for the performance prediction model. Third, the
evaluation will be extended to larger TPC-H scale factors (scale factor 10 and
scale factor 100) to characterize QueryPilot's behaviour on datasets where
MySQL's optimizer is more likely to make suboptimal decisions. Fourth, the
Python SDK will be extended to a transparent MySQL wire-protocol proxy,
allowing QueryPilot to intercept queries from any application without requiring
code changes.
